\par\medskip
\hypertarget{native-070-article053}{}\pdfbookmark[2]{53.}{article-053}
\begin{center}53.\end{center}
Die Lage eines beliebigen Punktes im Raum bestimmt man am bequemsten durch seine Abstände von drei Ebenen, die einander rechtwinklig schneiden. Wählen wir als eine dieser Ebenen die Ebene der Ekliptik und bezeichnen wir mit $z$ den Abstand des Himmelskörpers von ihr, der auf der Nordseite positiv und auf der Südseite negativ gezählt wird, so gilt offensichtlich
\[
z=r'\operatorname{tang}\beta=r\sin\beta=r\sin i\sin u.
\]
Die beiden übrigen Ebenen, die wir ebenfalls durch die Sonne gelegt denken, erzeugen auf der Himmelskugel Großkreise. Diese schneiden die Ekliptik rechtwinklig, sodass ihre Pole auf der Ekliptik selbst liegen und um $90^0$ voneinander entfernt sind. Den Pol einer Ebene, auf dessen Seite die Abstände positiv gezählt werden, nennen wir den \textit{positiven Pol}. Es seien also $N$ und $N+90^0$ die Längen der positiven Pole, und die Abstände von den entsprechenden Ebenen seien mit $x$ beziehungsweise $y$ bezeichnet. \SourceLeaf{071}Dann sieht man leicht, dass gilt
\[
\begin{aligned}
x&=r'\cos(\lambda-N)\\
&=r\cos\beta\cos(\lambda-\AscendingNode)\cos(N-\AscendingNode)\\
&\quad+r\cos\beta\sin(\lambda-\AscendingNode)\sin(N-\AscendingNode)\\[.4em]
y&=r'\sin(\lambda-N)\\
&=r\cos\beta\sin(\lambda-\AscendingNode)\cos(N-\AscendingNode)\\
&\quad-r\cos\beta\cos(\lambda-\AscendingNode)\sin(N-\AscendingNode),
\end{aligned}
\]
Diese Ausdrücke gehen über in
\[
\begin{aligned}
x&=r\cos(N-\AscendingNode)\cos u+r\cos i\sin(N-\AscendingNode)\sin u\\
y&=r\cos i\cos(N-\AscendingNode)\sin u-r\sin(N-\AscendingNode)\cos u.
\end{aligned}
\]
Legt man also den positiven Pol der zu $x$ gehörigen Ebene in den aufsteigenden Knoten selbst, sodass $N=\AscendingNode$ ist, so erhält man die einfachsten Ausdrücke für die Koordinaten $x$, $y$ und $z$
\[
\begin{aligned}
x&=r\cos u\\
y&=r\cos i\sin u\\
z&=r\sin i\sin u.
\end{aligned}
\]
Trifft diese Annahme nicht zu, so nehmen die oben angegebenen Formeln dennoch eine fast ebenso bequeme Form an, wenn man vier Hilfsgrößen $a$, $b$, $A$ und $B$ einführt, die so bestimmt sind, dass gilt\HistoricalNoteEight{Handschriftliche Bemerkung, wozu vgl. Seite 69 Zeile 7:]
\[
\begin{aligned}
-\cos\AscendingNode'&=a\sin A\\
\cos i'\sin\AscendingNode'&=a\cos A\\
-\sin\AscendingNode'&=b\sin B\\
-\cos i'\cos\AscendingNode'&=b\cos B.
\end{aligned}
\]
}
\[
\begin{aligned}
\cos(N-\AscendingNode)&=a\sin A\\
\cos i\sin(N-\AscendingNode)&=a\cos A\\
-\sin(N-\AscendingNode)&=b\sin B\\
\cos i\cos(N-\AscendingNode)&=b\cos B
\end{aligned}
\]
(siehe Art. 14, II). Dann gilt offenbar
\[
\begin{aligned}
x&=ra\sin(u+A)\\
y&=rb\sin(u+B)\\
z&=r\sin i\sin u.
\end{aligned}
\]
